\section{Justificación}
\label{sec:justificacion}

    El desarrollo de Colivi surge de una necesidad práctica vivida en primera persona durante estancias de intercambio universitario. Buscar habitación a cientos de kilómetros revela un mercado opaco. Los anuncios omiten detalles esenciales, los arrendadores tardan días en contestar y las condiciones económicas varían sobre la marcha sin dejar constancia escrita. El estudiante asume un desembolso económico sin garantías reales sobre el estado del piso ni la identidad de quien cobra la fianza.

    La convivencia posterior destapa nuevos roces cotidianos. Una vez formalizada la estancia, el grupo debe coordinar turnos de limpieza, compras compartidas y pagos de suministros. Confiar este control a pizarras domésticas o chats grupales suele terminar en disputas. La ausencia de un registro inmutable dificulta comprobar quién ha cumplido sus obligaciones y quién acumula retrasos continuados en los gastos.

    Las aplicaciones comerciales resuelven únicamente fragmentos aislados de este ciclo. Existen portales inmobiliarios centrados en captar inquilinos y aplicaciones móviles diseñadas para anotar gastos entre amigos, pero ninguna herramienta enlaza la reserva formal con la vida diaria dentro de la vivienda. Colivi conecta ambos momentos bajo un mismo sistema. La plataforma combina un backend en Spring Boot 3 con tipado estricto, una interfaz reactiva construida en React y un servidor auxiliar bajo el protocolo MCP (\emph{Model Context Protocol}) para consultar datos mediante modelos de lenguaje, cerrando la brecha entre la búsqueda inicial y el régimen diario del piso.
